\section{Uniform exponents for klt pluricanonical characters}
\label{sec:characters}

The rank bound now gives a uniform exponent for the scalar by which a
crepant birational transformation acts on a log volume form.  We work
over $\C$ throughout this section.  The transformation itself may have
infinite order.

\begin{theorem}[Uniform klt character exponent]\label{thm:characters}
For integers $n\geq0$ and $m\geq1$, there is an integer $L(n,m)>0$
with the following property.  Let $(X,\Delta)$ be an integral
projective klt $n$-fold pair, with $\Delta\geq0$ and
$m(K_X+\Delta)\sim0$ as an integral principal divisor.  Choose a
rational $m$-canonical form $\theta$ with
$\Div(\theta)=-m\Delta$.  For every B-birational self-map $f$ of
$(X,\Delta)$, write $f^*\theta=c\theta$, where $c\in\C^*$.
Then $c^{L(n,m)}=1$.
\end{theorem}

The proof separates the action into its rationally connected,
abelian, and nonabelian boundary-free factors.  The main issue is that
an arbitrary birational transformation need not preserve a given
product cover.  We first construct a cover admitting a lift, and then
recover the factor fields from their tangent directions.  This costs
only the exponent $n!$, independently of the degrees of the covers.

\subsection{A terminal model for the action}

Passing to a terminal pair makes the action an isomorphism in
codimension one.  We will then lift it through a finite cover of
the smooth locus.

\begin{lemma}\label{char:terminal-model}
In the setting of Theorem~\ref{thm:characters}, there is a projective
crepant birational morphism
\[
                 (V,B_V)\longrightarrow(X,\Delta)
\]
such that $V$ is $\Q$-factorial, $B_V\geq0$, and the pair $(V,B_V)$
is terminal.  Every B-birational self-map of $(X,\Delta)$ induces a
pseudo-automorphism of $V$ preserving $B_V$.
\end{lemma}

\begin{proof}
The extraction and $\Q$-factorialization theorem for klt pairs
\citep[Corollary~1.4.3]{BCHM2010} extracts the exceptional divisors
with discrepancy at most zero.  There are only finitely many such
divisors.  To recall why, take a log resolution with crepant
subboundary $\sum b_jD_j$, where every $b_j<1$.  For an exceptional
divisorial valuation $v$ over this resolution,
\[
 A_{\sum b_jD_j}(v)
 =A_{\sum D_j}(v)+\sum_j(1-b_j)v(D_j).
\]
The first term is a nonnegative integer.  If it is positive and the
centre is contained in the support, the displayed sum is strictly
greater than one.  If the generic point of the centre lies outside
the support, smoothness gives log discrepancy at least two.  Thus valuations of log discrepancy at most
one are lc places of the reduced simple normal crossings divisor.
Such places are toroidal; their positive integral weights are bounded
by $\sum_j(1-b_j)v(D_j)\leq1$.  Together with the finitely many
divisors on the resolution, this proves finiteness.  Extracting these
places gives the asserted terminal pair, and its boundary is
effective because each extracted discrepancy is nonpositive.

Pull $\theta$ back to $V$.  Its order at every prime divisor of $V$
is nonpositive, whereas its order at every exceptional divisor over
$V$ is strictly positive by terminality.  A transformation multiplying
$\theta$ by a nonzero constant cannot exchange these two kinds of
divisorial valuations.  Neither it nor its inverse therefore
contracts a prime divisor.  It is a pseudo-automorphism, and the
orders of $\theta$ show that it preserves the boundary coefficients.
\end{proof}

\subsection{Product covers and their factor fields}

We state the structural input in the form needed here.  A finite
surjective morphism of normal varieties is \emph{quasi-\'etale} if
it is \'etale in codimension one.  By purity it is \'etale over the
smooth locus of its target.  We write $\Omega^{[j]}$ for the
reflexive extension of $j$-forms from the smooth locus.

\begin{externalinput}[Product cover]\label{char:product-input}
Let $(V,B_V)$ be an integral projective klt pair with $B_V\geq0$,
$K_V$ rational Cartier, and $K_V+B_V\sim_{\Q}0$.  There is a finite
quasi-\'etale cover
\[
                  P=\prod_{i\in I}P_i\longrightarrow V
\]
whose positive-dimensional factors have the following properties.
There is at most one abelian factor and at most one rationally
connected factor $(F,H)$, where $(F,H)$ is klt, $H\geq0$, and
$K_F+H\sim_{\Q}0$; the entire pulled-back boundary is the pullback
of $H$, and is zero if that factor is absent.  All factors are
$\Q$-Gorenstein.  All other factors are
irreducible Calabi--Yau or
irreducible holomorphic symplectic varieties.  They, and every
connected normal finite quasi-\'etale cover of them, have canonical
singularities and Cartier trivial canonical class.  Their reflexive
form algebra is generated by a top form in the Calabi--Yau case and
by a symplectic form in the symplectic case.  In particular they have
no reflexive one-forms or vector fields.

The cover can be chosen so that every connected normal finite
quasi-\'etale cover of $F$ is rationally connected.  On every such
cover, reflexive one-forms and vector fields tangent to the
pulled-back boundary vanish.
\end{externalinput}

The product decomposition is proved in
\citet[Proposition~5.1]{LOGI}, using
the decomposition theorem of
\citet[Theorem~1.3]{MatsumuraWang2025} and the singular
Beauville--Bogomolov theorem
\citep[Definition~1.4 and Theorem~1.5]{HoringPeternell2019}.  The statements about
further covers and logarithmic vector fields are respectively
Lemmas~5.3 and~5.2 of \citet{LOGI}.  The abelian factor includes
all one-dimensional boundary-free factors.  A connected normal
finite quasi-\'etale cover of an abelian variety is \'etale, and is
again abelian after choosing an origin.

The following argument adapts the factor-direction construction in
\citet[Proposition~5.4]{LOGI}.  That proposition treats finite regular actions;
we give the argument needed for arbitrary pseudo-automorphisms.

\begin{lemma}[Lifting and separating a birational action]
\label{char:factor-actions}
Let $(V,B_V)$ be as in Lemma~\ref{char:terminal-model}, of dimension
$n>0$.  Choose a product cover $P=\prod_{i\in I}P_i\to V$ as in
Input~\ref{char:product-input}, and let $f$ be a B-birational
self-map with $f^*\theta=c\theta$.  There are normal integral projective
varieties $R,R_i$, finite quasi-\'etale maps
\[
 R\xrightarrow{\ \psi\ }\prod_{i\in I}R_i
     \longrightarrow\prod_{i\in I}P_i\longrightarrow V,
\]
and birational self-maps $h_i$ of $R_i$ with the following properties.
Each $R_i\to P_i$ is a finite quasi-\'etale cover.  Write $H_i$ for
its boundary, zero except on the rationally connected factor, and
put $p_i=m$ on that factor and $p_i=1$ on the others.  There are
rational $p_i$-canonical forms $\theta_i$ with
\[
 \Div(\theta_i)=-p_iH_i,\qquad h_i^*\theta_i=d_i\theta_i
 \quad(d_i\in\C^*),
\]
and
\begin{equation}\label{char:scalar-product}
                       c^{n!}=\prod_{i\in I}d_i^{m/p_i}.
\end{equation}
In particular each $h_i$ is B-birational for $(R_i,H_i)$.
\end{lemma}

\begin{proof}
\emph{A cover admitting a lift.}
Put $U=V_{\mathrm{sm}}$.
The connected \'etale cover $P_U\to U$ corresponds to a finite-index
subgroup of $\pi_1^{\mathrm{top}}(U)$.  This fundamental group is
finitely generated, so it has only finitely many subgroups of any
fixed index.  Their intersection is a finite-index characteristic
subgroup contained in the subgroup defining $P_U$.  Riemann
existence gives its connected algebraic \'etale cover.  Normalize
$V$ in its function field to obtain a finite quasi-\'etale cover
$R\to V$ factoring through $P$.

The pseudo-automorphism $f$ identifies two open subsets of $U$ with
complements of codimension at least two.  Removing such subsets
does not change the fundamental group of a complex manifold, by
transversality.  Characteristicity therefore lifts this isomorphism
to the chosen cover.  Riemann existence makes the lift algebraic,
and it induces a pseudo-automorphism $g$ of $R$.  If $H_R$ and
$\theta_R$ denote the pulled-back boundary and form, then
\[
                 g^*\theta_R=c\theta_R,
\]
and $g$ preserves $H_R$ in codimension one.  No degree bound for
$R\to V$ is required.

\emph{Preserving the factor directions.}
Over the product of the smooth loci of the $P_i$, the map $R\to P$
is \'etale.  The tangent directions of the factors give a splitting
there, whose reflexive extension is
\[
                         \mathcal T_R=\bigoplus_{i\in I}\mathcal E_i.
\]
Let $\mathcal T_R(-\log\Supp H_R)$ be the reflexive sheaf of
derivations preserving the reduced ideal of every boundary component,
a condition tested at codimension-one points.  Consider the
finite-dimensional $\C$-algebra
\[
 \mathcal A=\{u\in H^0(R,\End(\mathcal T_R)):
 u(\mathcal T_R(-\log\Supp H_R))
       \subseteq\mathcal T_R(-\log\Supp H_R)\}.
\]
Conjugation by $g$ acts on $\mathcal A$: both the tangent sheaf and
the logarithmic condition are determined on a big open set.

Every member of $\mathcal A$ has zero entries between different
$\mathcal E_i$.  To see this, fix general smooth points in the
complementary factors and slice in one factor direction.  Each
connected component over that factor's smooth locus extends by
normalization to a connected normal finite quasi-\'etale cover of
the factor.  The omitted set has codimension at least two, since a
divisorial point of a finite cover lies over a divisorial, hence
smooth, point of the normal factor.  The ambient cover is \'etale
on this open slice, so restricted entries are regular and extend
reflexively across its complement.  On the slice its tangent
directions are $\mathcal E_i$, while complementary summands are
trivial.  An entry
in either direction involving a nonabelian boundary-free factor
therefore gives reflexive one-forms or vector fields on its cover,
and these vanish by Input~\ref{char:product-input}.  For the only
remaining case, an abelian and a rationally connected factor, slice
in the rationally connected direction.  An entry towards the
abelian summand gives one-forms; an entry in the reverse direction
gives vector fields tangent to the boundary: that boundary is
pulled back from the rationally connected factor, so the
ideal-preservation condition restricts at every divisorial point
of the slice.  These vanish as well.  The slices
cover a dense open of $R$, proving the assertion.

Consequently every block projection is a central idempotent of
$\mathcal A$.  The primitive central idempotents give nonzero
direct summands of the rank-$n$ tangent sheaf, so there are at most
$n$ of them.  Conjugation by $g$ permutes these idempotents.
Therefore
\[
                              h=g^{n!}
\]
fixes every primitive central idempotent.  Every central idempotent
is a sum of primitive central idempotents, so $h$ fixes every block
projection and preserves each factor direction.

\emph{Recovering the factor fields.}
Let $K_i$ be the relative algebraic closure of $\C(P_i)$ in
$\C(R)$, and let $R_i$ be the normalization of $P_i$ in $K_i$.
Equivalently, $R\to R_i\to P_i$ is the Stein factorization of the
projection.  At the generic point of $R$, the directions
complementary to $i$ span
$\operatorname{Der}_{\C(P_i)}\C(R)$.  In characteristic zero their
common field of constants is exactly $K_i$.  Preservation of these
directions thus gives a birational self-map $h_i$ of $R_i$.

The map $\psi:R\to\prod_iR_i$ is proper and has finite fibres,
because its composite to $P$ is finite.  It is finite and, by
dimension, surjective.  The maps $R_i\to P_i$ are quasi-\'etale.
Indeed, since $\C(P)/\C(P_i)$ is a regular extension, the product
$R_i\times\prod_{j\ne i}P_j$ is an intermediate normal finite
cover of $P$.  Ramification over a prime of $P_i$ would persist
in this product and then in a prolongation to $R$, contradicting
quasi-\'etaleness of $R\to P$.  The same ramification calculation
in the tower proves that $\psi$ is quasi-\'etale.  In particular the
pulled-back pairs $(R_i,H_i)$ are klt.  The relation
between $h$ and the $h_i$ is rational equivariance of $\psi$; no
regularity assertion about the $h_i$ is needed.

\emph{Comparing the forms.}
On each boundary-free $R_i$, choose a generator $\theta_i$ of its
trivial canonical sheaf.  For the rationally connected factor,
restrict the degree-$m$ trivialization pulled back from $V$ to that
factor of $P$, using smooth complementary points and the external
product of canonical lines.  The resulting form has divisor $-mH$
on that factor; pull it back to $R_i$.  This constructs the forms
of degree $p_i$ in the statement.  Their product, in any fixed
order, satisfies
\begin{equation}\label{char:form-product}
 \theta_R=a\,\psi^*\left(\bigwedge_i\theta_i^{\otimes m/p_i}\right)
 \qquad(a\in\C^*).
\end{equation}
Here the wedge denotes external product of pluricanonical lines.
Indeed both sides have divisor $-mH_R$, so their ratio is a global
unit on the normal projective integral variety $R$.

Write $u_i=h_i^*\theta_i/\theta_i\in\C(R_i)^*$.  Pulling
\eqref{char:form-product} back by $h$ shows that the pullback of
$\prod_i u_i^{m/p_i}$ is the constant $c^{n!}$.  Since $\psi$ is
dominant, the product is constant on $\prod_iR_i$.  Fixing general
points in all complementary factors shows that every
$u_i^{m/p_i}$, and hence every $u_i$, is constant.  Denote it by
$d_i$.  This proves \eqref{char:scalar-product}; proportionality of
the forms proves that each $h_i$ is B-birational.
\end{proof}

We have reduced the character on $V$ to characters on finite covers
of the individual factors.  Their dimensions are at most $n$, and the
integer $m$ has not changed.  We next bound these factor characters, using
bounded families on the rationally connected factor and the rank
bound on the other nonabelian factors.

\subsection{Characters on the factors}

\begin{lemma}[Rationally connected factors]\label{char:rc-bound}
For $r,m\geq1$ there is an integer $E_{\mathrm{rc}}(r,m)>0$ that
kills the degree-$m$ B-birational character of every integral
projective rationally connected klt pair $(F,H)$ of dimension $r$
with $H\geq0$ and $m(K_F+H)\sim0$ integral principal.
\end{lemma}

\begin{proof}
By \citet[Theorem~1.2]{HanJiang}, the underlying varieties are
bounded modulo flops: they are isomorphic in codimension one to
normal members $Y$ of a bounded family.  The theorem applies because
$-m(K_F+H)$ is Cartier and nef.  In dimension one one can instead
use $F\simeq\PP^1$ directly.

Let $B_Y$ be the strict transform of $H$.  A degree-$m$ log
trivializing form has divisor $-mB_Y$ on $Y$, by the codimension-one
identification.  Thus $m(K_Y+B_Y)\sim0$ is integral principal, and
the map is B-birational by the form criterion.  In particular the
transformed pair is klt.  Choose very ample divisors $A$ in bounded
projective embeddings of these $Y$, with $A^r$ uniformly bounded.
Since $mB_Y$ is integral,
\begin{equation}\label{char:boundary-degree}
 (\Supp B_Y)A^{r-1}
 \leq mB_YA^{r-1}=-mK_YA^{r-1}
 \leq m\bigl(2+(r-1)A^r\bigr).
\end{equation}
For the last inequality, take a general complete-intersection curve
of $r-1$ members of $|A|$.  It is smooth and avoids the singular
locus of the normal variety $Y$, and adjunction gives
$-K_YA^{r-1}=2-2g+(r-1)A^r$.

Reduced subvarieties of bounded degree in bounded projective
embeddings form a bounded family.  Moreover, the nonzero
coefficients of $B_Y$ lie in
$\{1/m,\ldots,(m-1)/m\}$, since the pair is klt.  Thus the pairs
are log bounded modulo B-birational contractions.  (A birational
map that is an isomorphism in codimension one is a birational
contraction in this terminology.)

\citet[Theorem~3.2]{JiangLiu}, which applies in every dimension,
now gives integers $k,N>0$ depending only on $r,m$ such that
$k(K_F+H)\sim0$ and the image of the degree-$k$ B-representation
has order at most $N$.  If a map acts by $d$ on a degree-$m$
generator and by $e$ on a degree-$k$ generator, then comparison in
degree $km$ gives $d^k=e^m$.  Since $e^{N!}=1$, we obtain
$d^{kN!}=1$.  We may take $E_{\mathrm{rc}}(r,m)=kN!$.
\end{proof}

\begin{lemma}[Boundary-free factors]\label{char:hodge-bound}
For each $r\geq1$ there is an integer $E_0(r)>0$ with the following
property.  Let $Y$ be a connected normal finite quasi-\'etale cover
of an $r$-dimensional abelian, irreducible Calabi--Yau, or irreducible
holomorphic symplectic factor of Input~\ref{char:product-input}.
Every birational self-map of $Y$ acts on its canonical volume form
by a scalar killed by $E_0(r)$.
\end{lemma}

\begin{proof}
Let $U$ be a smooth projective resolution of $Y$.  The canonical
singularities and trivial canonical class give
$h^{r,0}(U)=1$.  For a nonabelian factor,
\citet[Lemma~5.2]{UPI}
shows that $q(U)=0$ and that every positive-dimensional rational
image of smaller dimension has rationally connected smooth
resolution.  Its hypotheses hold for $Y$ itself, using the identity
cover, because the defining form-algebra properties persist under
all further connected normal finite quasi-\'etale covers.
Take a crepant projective $\Q$-factorial terminalization of $Y$.
It still has principal canonical divisor, and the two properties
just stated are birational invariants.  Theorem~\ref{thm:rank}
therefore bounds $\dim_{\Q}T(U)$ in terms of $r$.  Nonabelian
factors here have dimension at least two.  For an abelian factor,
take $U=Y$ and use instead
\[
                  \dim_{\Q}T(U)\leq\dim_{\Q}H^r(U,\Q)
                                      =\binom{2r}{r}.
\]
Thus there is a common dimension bound $b(r)$ in either case.

Let $p,q:W\to U$ be a smooth projective common resolution of the
induced birational self-map of $U$.  By Lemma~\ref{hodge:basic},
birational invariance of the Hodge structure generated by the top
form gives
\[
                     T(W)=p^*T(U)=q^*T(U).
\]
Consequently $A=p_*q^*$ preserves $T(U)$ and the cup-product
pairing there; its inverse is $q_*p^*$.  Pullback and pushforward
are integral on cohomology modulo torsion, so $A$ preserves the full
lattice
\[
        T(U)\cap\bigl(H^r(U,\Z)/\text{torsion}\bigr).
\]
Moreover, $T(U)$ is primitive for every rational ample class:
the kernel of cup product by that class is a rational Hodge
substructure containing $H^{r,0}(U)$.  The Hodge--Riemann relations
therefore turn the cup pairing and Hodge decomposition on $T(U)$
into a positive-definite Hermitian form, preserved by $A$.

The scalar $d$ on the canonical volume form is an eigenvalue of
this integral operator.  It is an algebraic integer, and all its
conjugates occur among the eigenvalues of $A$, hence have modulus
one.  Kronecker's theorem makes $d$ a root of unity.  If its order
is $e$, the degree of its cyclotomic polynomial is Euler's totient
$\varphi(e)$, so
\[
                         \varphi(e)\leq b(r).
\]
There are only finitely many positive integers with this property.
Their least common multiple is a permissible choice of $E_0(r)$.
\end{proof}

\begin{proof}[Proof of Theorem~\ref{thm:characters}]
For $n=0$ the pair is a point and the character is trivial, so take
$L(0,m)=1$.  Suppose $n>0$.  Lemmas~\ref{char:terminal-model} and
\ref{char:factor-actions} express $c^{n!}$ as in
\eqref{char:scalar-product}.  The rationally connected factor
satisfies the hypotheses of Lemma~\ref{char:rc-bound}; in particular
it remains rationally connected by the cover-stable choice in
Input~\ref{char:product-input}.  Every other factor satisfies
Lemma~\ref{char:hodge-bound}.  Define
\[
 E(n,m)=\operatorname{lcm}_{1\leq r\leq n}
          \{E_{\mathrm{rc}}(r,m),E_0(r)\},
 \qquad L(n,m)=n!E(n,m).
\]
Each $d_i^{E(n,m)}=1$, and \eqref{char:scalar-product} gives
$c^{L(n,m)}=1$.  Every bound depends only on the factor dimensions
and $m$, and therefore only on $n,m$.
\end{proof}
